\section{Related Work}
\label{sec:related}

\subsection{LLM Uncertainty Quantification}

\paragraph{Semantic entropy.}
Kuhn et al.~\cite{kuhn2023semantic} introduce semantic entropy as a method for
detecting hallucinations in free-form text generation. Multiple responses to
the same prompt are generated and clustered by semantic equivalence (via
natural language inference); high entropy across clusters indicates uncertainty.
The method is principled and effective for factual question answering, where
model uncertainty produces cross-generation variance. As we show in
Section~\ref{sec:comparison}, this variance is not present for code generation
API decisions: models generate the same (potentially wrong) API call across
all samples.

\paragraph{UQLM.}
The UQLM library~\cite{uqlm2024} provides open-source implementations of
multiple uncertainty scoring methods for LLM outputs, including semantic
negentropy, semantic clustering, and conformal prediction wrappers. It is the
closest published tool to our work and serves as our primary baseline.
Our direct comparison (Section~\ref{sec:comparison}) shows 5\% detection rate
(1 of 20 high-uncertainty prompts) compared to 85\% for $\eps$.

\paragraph{Conformal prediction for LLMs.}
Quach et al.~\cite{quach2023conformal} and Angelopoulos et al.~\cite{angelopoulos2022conformal}
apply conformal prediction to construct prediction sets with coverage
guarantees over LLM outputs. The approach is statistically rigorous but
produces set-valued predictions over outputs rather than within-generation
token-level signals. It also requires a calibration dataset and does not
provide the localization (which token, which line, which function) that
makes $\eps$ actionable for code review.

\paragraph{Token-level calibration.}
Kadavath et al.~\cite{kadavath2022language} show that LLMs are reasonably
well-calibrated at the token level: stated confidence correlates with accuracy.
Varshney et al.~\cite{varshney2023stitch} use token-level log-probabilities
for selective generation, abstaining when confidence is low. Our work extends
this line by identifying \emph{which} tokens carry meaningful uncertainty in
code --- specifically, by filtering out cosmetic naming decisions that dominate
naive entropy measurements. Uncertainty propagation through multi-stage ML
pipelines is well-studied in adjacent domains~\cite{jungo2020analyzing}; the
$\eps$ primitive is designed to support analogous propagation in code
generation pipelines, deferred to future work.

\paragraph{Code generation uncertainty and self-consistency.}
Concurrent approaches to code generation quality estimation include
test-based ranking (CodeT~\cite{chen2022codet} generates tests and ranks
candidate solutions by pass rate) and repeated-sampling self-consistency
(AlphaCode-style approaches use pass@$k$ as an implicit uncertainty proxy).
These methods are inter-generational by construction: they require multiple
independent samples and aggregate across them. They are therefore subject to
the same structural limitation as UQLM for the problem we study --- a model
that consistently generates the same API version across samples will produce
consistent test outcomes regardless of whether the API version is correct.
Our approach is complementary: it operates within a single generation and is
zero-overhead.

\subsection{Code Generation Quality}

Recent work on evaluating LLM code generation quality focuses primarily on
functional correctness via execution-based benchmarks (HumanEval~\cite{chen2021evaluating},
MBPP~\cite{austin2021program}, SWE-bench~\cite{jimenez2024swebench}). These
benchmarks measure whether generated code passes test cases, not whether the
model was confident in its decisions. A model that passes all tests with low
$\eps$ is reliable; a model that passes all tests with high $\eps$ on specific
API decisions may be one dependency update away from failure.

To our knowledge, no prior work treats within-generation token-level
uncertainty as a first-class signal for code generation risk. The closest work
is in automated program repair, where uncertainty has been used to rank repair
candidates~\cite{ye2022selfapr}, but this operates post-generation and does
not address the intra-generation measurement problem we identify.
